\chapter*{Reading Paths}
\addcontentsline{toc}{chapter}{Reading Paths}
\index{reading paths}

The book is designed for selective reading. Chapters state their own evidence boundaries, while
Appendix~\ref{ch:verification-tiers} defines the shared vocabulary. The routes below avoid
repeating material and identify a useful first destination for common readers.

\begin{description}[style=nextline,leftmargin=3.2cm,labelwidth=2.8cm]
  \item[New CNA user]
    Read Chapters~\ref{ch:what-is-cna}, \ref{ch:configuring-building}, and
    \ref{ch:first-game}. Continue with the game lifecycle (Chapter~\ref{ch:game-lifecycle}),
    SpriteBatch (Chapter~\ref{ch:spritebatch}), and content resolution
    (Chapter~\ref{ch:content-and-assets}). Chapter~\ref{ch:renderer-selection-identity} explains renderer
    choice; Chapter~\ref{ch:roadmap} lists material gaps at the edition pin.

  \item[XNA or FNA porter]
    Start with the compatibility boundary in Chapters~\ref{ch:cna-xna}
    and~\ref{ch:design-philosophy}, then use the migration workflow in
    Chapter~\ref{ch:migration-guide}. Consult Chapters~\ref{ch:content-and-assets}--
    \ref{ch:cnj-format} for assets, Chapter~\ref{ch:shader-fx-gap} for effects and custom
    shaders, and Appendix~\ref{ch:cnaext-catalog} when strict XNA compatibility matters.

  \item[Renderer implementer]
    Read Chapters~\ref{ch:graphicsdevice}, \ref{ch:backend-architecture},
    \ref{ch:renderer-selection-identity}, and~\ref{ch:vertex-streams-capabilities}. Appendix~\ref{ch:feature-matrix}
    summarizes the family-level contract; Chapters~\ref{ch:hostile-environments}
    and~\ref{ch:oracles} explain engagement and oracle requirements.

  \item[Runtime renderer user]
    Read Chapters~\ref{ch:configuring-building}, \ref{ch:renderer-selection-identity}, and
    \ref{ch:backend-architecture}. The selection chapter separates compiled families, configured
    default, runtime preference, fallback and the first-device latch; Appendix~\ref{ch:feature-matrix}
    supplies the family inventory and platform gates.

  \item[Native C consumer]
    Begin with Appendix~\ref{ch:native-c-api}, then use Chapters~\ref{ch:configuring-building}
    and~\ref{ch:module-monorepo} for build and link topology. Read Chapter~\ref{ch:ci-reality}
    before interpreting the C API workflow and release-gate inventory as executed evidence.

  \item[Framework contributor]
    Begin with the lifecycle and module boundaries in Chapters~\ref{ch:game-lifecycle},
    \ref{ch:module-monorepo}, and~\ref{ch:boundaries-enforcement}. Chapters~\ref{ch:cnatests-architecture}
    and~\ref{ch:discipline-as-test} cover test registration and mechanical policy gates.

  \item[Content or glTF contributor]
    Read Chapters~\ref{ch:content-and-assets}--\ref{ch:content-robustness}, then
    Chapters~\ref{ch:gltf-import-core}--\ref{ch:gltf-conformance}. Appendix~\ref{ch:gltf-evidence-matrix}
    keeps importer, runtime, renderer, and oracle evidence in separate columns.

  \item[Verification contributor]
    Read Parts~\ref{part:testing-verification} and Appendix~\ref{ch:verification-tiers} first.
    Chapter~\ref{ch:counting-discipline} defines reproducible counts; Chapter~\ref{ch:oracles}
    distinguishes authority and tolerance; Chapter~\ref{ch:ci-reality} records what automation
    did and did not exercise at the pin.
\end{description}

\section*{The shortest practical route}

For a first program, the minimum sequence is: select a renderer when configuring CNA, build CNA
as the repository documents, derive a class from \cnaclass{Game}, let
\cnaclass{GraphicsDeviceManager} create the graphics device, load assets through
\cnaclass{ContentManager}, and draw through an XNA-shaped API such as \cnaclass{SpriteBatch}.
Chapter~\ref{ch:first-game} contains the complete small program and commands. Portable code stays
within the \texttt{Microsoft::Xna::Framework} surface and treats \texttt{CNAEXT}-marked APIs as
explicit extensions.

When a result differs by renderer, first identify the public selector, implementation family,
native or translation API, and whether the intended route engaged. Then reduce the case to a
state transition, readback, or pixel oracle. A successful process exit alone cannot localize a
renderer-specific failure.
